% title: Ions
An atom is normally electrically neutral: it carries the same number of
negatively charged electrons as positively charged protons, so it holds no net
electrical charge. An \textbf{ion} is an atom, or a compound, that has gained
or lost one or more electrons and therefore carries a net electrical
charge. That charge governs how ions bond, how they move, and how they carry
electrical signals. Ions come up very often in our UCSCICOG11 course, so it is
worth making sure you understand what they are and what their role is in neurophysiology.

\subsection{Cations and anions}
When an atom loses one or more electrons it is left with more protons than
electrons and so becomes positively charged; a positive ion is called a
\textbf{cation}. When an atom gains electrons it becomes negatively charged and
is called an \textbf{anion}. The size of the charge is written as a superscript:
a sodium atom that has lost one electron is Na$^{+}$, a calcium atom that has
lost two is Ca$^{2+}$, and a chlorine atom that has gained one electron is
Cl$^{-}$ (in that form it is called chloride).

Four ions do most of the electrical work in the body, including the nervous system:
\begin{itemize}
  \item \textbf{Sodium} (Na$^{+}$): the main cation of the fluid
    \emph{outside} cells.
  \item \textbf{Potassium} (K$^{+}$): the main cation \emph{inside} cells.
  \item \textbf{Chloride} (Cl$^{-}$): the main anion of the fluid outside
    cells.
  \item \textbf{Calcium} (Ca$^{2+}$): held at a very low concentration inside
    cells, so even a small inflow acts as a strong signal.
\end{itemize}

\subsection{Extracellular and intracellular ions}
Every cell is wrapped in a thin \textbf{cell membrane} that separates two watery
compartments: the \textbf{intracellular fluid} (the cytoplasm, inside the cell)
and the \textbf{extracellular fluid} (the surroundings, outside the cell). These
two compartments do not hold the same concentrations of ions. Using energy, the
cell runs molecular pumps in its membrane, most importantly the
\textbf{sodium--potassium pump} that pushes Na$^{+}$ out of the cell and
K$^{+}$ in. As a result, potassium is concentrated inside the cell while sodium,
chloride, and calcium are concentrated outside it. These differences persist
because ions, being charged, cannot slip across the membrane on their own; they
can only pass through dedicated protein pores called \textbf{ion channels}, and
only when those channels are open. More about this in our course.

\subsection{How ions make electricity}
Separating positive from negative charge across a barrier is exactly what a
battery does. Because the intracellular and extracellular fluids hold different
amounts of charged ions, a small `voltage' (a difference in electrical potential)
exists across the cell membrane. This is the \textbf{membrane potential}. In
a resting neuron the inside of the cell sits at about $-70$~millivolts relative
to the outside. In our course, you'll learn how this works. Very briefly: when
ion channels open, ions flow across the membrane depending on the
concentration and electrical gradients. A flow of charge is an \textbf{electric current},
so opening and closing ion channels lets a cell change its membrane potential
quickly. This is the physical basis of \textbf{bioelectricity}. The \textbf{electrical
potentials} that neurons use to communicate, and the ion
channels that produce them, are an important foundational topic in the first part
of the course.
